\section{FORMAL RESTATEMENT}
\textbf{Erd\H{o}s \#296; a complete greedy-packing deduction, 2026-09-24.}
Let $k(N)$ be the largest number of pairwise disjoint subsets $A_1,\ldots,A_k\subseteq[1,N]$ satisfying $\sum_{n\in A_i}1/n=1$ for each $i$. Estimate $k(N)$ and decide whether $k(N)=o(\log N)$.

\section{QUICK LITERATURE AND CONTEXT CHECK}
The proposed little-oh estimate is false: $k(N)\sim\log N$. The problem page credits Hunter and Sawhney's observation that Bloom's reciprocal-mass theorem implies this by a greedy argument. We reproduce that deduction, using the stronger reciprocal-mass threshold of Liu and Sawhney as an explicit external input. The conclusion is a known result.

\section{OBJECTIVE AND ATTACK PLAN}
Every selected set consumes exactly one unit of the total harmonic mass. Continue selecting sets until the unused denominators contain no further unit-sum subset. A theorem bounding the reciprocal mass of such a remainder controls the total loss.

\section{WORK}
Write $R(E)=\sum_{n\in E}1/n$ and $H_N=R([1,N])$. Disjointness immediately gives
\[
k(N)=\sum_{i=1}^{k(N)}R(A_i)
=R\!\left(\bigcup_i A_i\right)\leq H_N.
\tag{1}
\]
In particular $k(N)\leq\lfloor H_N\rfloor$.

\textbf{Declared external input.} Liu and Sawhney's Theorem 1.1 says that for each fixed $\varepsilon>0$ and all sufficiently large $N$, every $E\subseteq[1,N]$ satisfying
\[
R(E)\geq(\log N)^{4/5+\varepsilon}
\tag{2}
\]
contains a subset with reciprocal sum exactly $1$. Distinctness is built into the use of subsets. The Fourier-analytic proof of this theorem is not reconstructed here.

Set $\varepsilon=1/10$. Starting with $[1,N]$, repeatedly remove a subset of reciprocal sum $1$ whenever one exists. Each removal deletes at least one integer, so the procedure stops after a finite number $q$ of steps. Its output consists of disjoint sets $A_1,\ldots,A_q$ and a remainder $U$ containing no unit-sum subset. By the contrapositive of (2),
\[
R(U)<(\log N)^{9/10}
\]
for sufficiently large $N$. The exact mass identity is
\[
H_N=q+R(U).
\tag{3}
\]
Therefore
\[
H_N-(\log N)^{9/10}<q\leq k(N)\leq H_N.
\tag{4}
\]
The procedure need not find an optimal packing: every maximal packing already satisfies the lower bound in (4). Using $H_N=\log N+\gamma+O(1/N)$ yields
\[
k(N)=\log N+O((\log N)^{9/10}),\qquad
\frac{k(N)}{\log N}\longrightarrow1.
\tag{5}
\]
More generally, any threshold $T(N)=o(\log N)$ in place of (2) would give $H_N-T(N)<k(N)\leq H_N$ by the same proof. This explains why Bloom's earlier estimate also suffices for the asymptotic.

\section{VERIFICATION AND SCOPE}
The remainder $U$ may be arbitrary and sparse; the external theorem is stated for arbitrary subsets of $[1,N]$, so it still applies after all the removals. Pairwise disjointness, rather than any independence assumption, gives the exact identity (3). The proof answers the main question about sums equal to $1$; the separate editorial question about many sets with an unspecified common sum is not asserted here.

\section{SOURCES}
\begin{itemize}
\item Current problem and the Hunter--Sawhney deduction: \url{https://www.erdosproblems.com/296}.
\item T. F. Bloom, \emph{On a density conjecture about unit fractions}, Theorem 3: \url{https://arxiv.org/html/2112.03726}.
\item Y. P. Liu and M. Sawhney, \emph{On further questions regarding unit fractions}, Theorem 1.1: \url{https://arxiv.org/html/2404.07113}.
\end{itemize}
\section{CONCLUSION}
\textbf{Attempt status: solved, using the stated reciprocal-mass theorem.}
In fact $k(N)\sim\log N$, so the proposed $o(\log N)$ bound is false.
\textbf{Completion rate estimate: 100\% of the main question, relative to the external input.}
